% Final Merits Determination
% EB-1A Petition - Vivek Kumar Mishra
% Drafted in accordance with Kazarian v. USCIS and AAO precedent

\subsection{Final Merits Determination: Sustained National Acclaim and Standing at the Top of the Field}
\label{sec:finalmerits}

The regulatory criteria establish the evidentiary baseline. The determination of whether the beneficiary has achieved sustained national or international acclaim and ranks among the small percentage at the very top of the field requires examination of the totality of the evidence.

\subsubsection{Nature of the Field}

The beneficiary's field is the development and deployment of artificial intelligence and machine learning systems for critical national infrastructure. This field is measured not by publication volume alone, but by institutional reliance: whether government agencies, systemically important financial institutions, and academic programs adopt and depend upon the beneficiary's methodologies and systems.

In this field, impact is demonstrated through federal archival of technical methodologies, cross-state adoption by government transportation agencies, deployment of production systems serving tens of millions of Americans, and selection to evaluate cutting-edge research at premier international conferences. These forms of validation reflect that the work has moved beyond academic contribution to become infrastructure upon which institutions rely.

\subsubsection{Independent Reliance and Field-Level Influence}

The beneficiary's work has been adopted by entities with no institutional connection to his employment or degree-granting institutions. The U.S. Department of Transportation's Bureau of Transportation Statistics has permanently archived two of the beneficiary's machine learning methodologies in the National Transportation Library (ROSA P Records dot/61210 and dot/59030). Federal agencies archive only those technical contributions deemed authoritative for long-term reference by transportation professionals and policymakers.

The Florida Department of Transportation subsequently adopted these methodologies in federally-funded research programs conducted at FAMU-FSU College of Engineering, as attested by Dr. Maxim A. Dulebenets, the principal investigator. This represents cross-state transfer: methodologies developed for California infrastructure challenges were determined sufficiently robust and generalizable for application to Florida's transportation systems. Dr. Dulebenets states that the beneficiary's work ``shaped aspects of our freight network modeling and equity-focused analyses in Florida.''

The beneficiary's employment at California State University, Long Beach provides further evidence of reliance. Dr. Shailesh Chandra, Assistant Professor in the Department of Civil Engineering and Construction Engineering Management, formally recruited the beneficiary in February 2020 specifically for ``programming skills in Python'' to develop machine learning frameworks for federally-funded intelligent transportation systems research. This cross-departmental recruitment—a Computer Science graduate student employed by Civil Engineering—demonstrates that the beneficiary possessed technical capabilities the Civil Engineering department itself lacked. The appointment letter explicitly identifies the research as funded by the U.S. Department of Transportation through the Mineta Transportation Institute, and the extension letter in September 2021 assigned development of ``Python code for keyword indexing and advanced text data processing,'' indicating sustained technical leadership over a two-year period on increasingly complex tasks.

At JPMorgan Chase \& Co., the beneficiary served as a founding engineer on the Chase Shopping Hub platform, which processes over 5 million customer requests monthly with 99.9\% reliability. Corey Graham, Executive Director of Software Engineering with 23 years of industry experience, attests that the beneficiary ``designed APIs that are now part of a mission-critical system'' and that his role was ``critical'' to the platform serving 80+ million Chase cardmembers. Sumit Bhatnagar, Vice President of Engineering, independently corroborates that the beneficiary ``took full ownership of the backend architecture'' and introduced cost-control solutions resulting in approximately \$100,000 in savings and performance improvements of 43\%.

The beneficiary has been selected to evaluate research submissions for the FM-Wild Workshop at the International Conference on Learning Representations (ICLR) 2025, the Position Paper Track at the Conference on Neural Information Processing Systems (NeurIPS) 2025, and the AIWILD Workshop at the International Conference on Machine Learning (ICML) 2026. ICLR, NeurIPS, and ICML are the premier venues for machine learning research globally, with main-track acceptance rates of 32\% and 24.52\% respectively for ICLR and NeurIPS. Reviewer selection for affiliated workshops and specialized tracks at these conferences is based on demonstrated domain expertise. The beneficiary completed reviews for multiple papers, including providing technical evaluation of autonomous multi-agent systems, AI applications to energy forecasting, and agentic AI safety. The IEEE, the world's largest technical professional organization, selected the beneficiary to serve on the Senior Member Application Virtual Review Panel in June 2025, a role requiring evaluation of whether applicants have demonstrated the ``significant performance'' required for Senior Member status.

\subsubsection{Evidence of Sustained Acclaim}

The record demonstrates recognition spanning multiple years and originating from independent institutions:

\begin{itemize}
\item 2018: Commanded salary 83\% above market average at Oracle Financial Services Software Limited
\item 2020--2021: Recruited by CSULB Civil Engineering department for specialized AI/ML capabilities; research outputs subsequently archived by U.S. DOT
\item 2021--present: Critical role at JPMorgan Chase, a Systemically Important Financial Institution
\item 2025--2026: Selected as reviewer for the ICLR 2025 FM-Wild Workshop, the NeurIPS 2025 Position Paper Track, and the ICML 2026 AIWILD Workshop; elevated to IEEE Senior Member; technical book adopted by government engineering institution (Supaul College of Engineering) for formal curriculum integration
\end{itemize}

These recognitions are not clustered in a single period but reflect continuous validation by unrelated entities over seven years.

\subsubsection{Distinction From the General Population of Practitioners}

The general population of AI/ML practitioners does not command federal archival of their methodologies in official U.S. government repositories. The general population does not serve as reviewers within the research ecosystems of ICLR, NeurIPS, and ICML---conferences with H5-indices of 362 and main-track submission volumes exceeding 20,000 papers, whose affiliated workshops and specialized tracks require demonstrated subject-matter expertise for reviewer selection. The general population does not hold critical roles at Fortune 13 systemically important financial institutions where platform failures would affect 84 million American consumers.

The cross-institutional pattern---federal government archival, state transportation agency adoption, Fortune 13 financial institution reliance, and selection to evaluate research within the ecosystems of premier AI/ML conferences---distinguishes the beneficiary from practitioners who may possess technical skills but whose work has not achieved field-level adoption and institutional dependence.

\subsubsection{National Importance and Benefit}

The beneficiary's work directly supports U.S. financial infrastructure stability and federal transportation priorities. JPMorgan Chase is classified as a Systemically Important Financial Institution under the Dodd-Frank Act; disruption to its systems would have cascading effects on the national economy. The beneficiary's contributions to intelligent transportation systems align with the Biden Administration's Justice40 Initiative and Infrastructure Investment and Jobs Act priorities for equitable, data-driven transportation planning.

Executive Order 14110 on Safe, Secure, and Trustworthy Artificial Intelligence (October 30, 2023) identifies AI development as critical to national security and economic competitiveness. The beneficiary's demonstrated capacity to develop production AI systems at scale, conduct peer-reviewed AI research adopted by government agencies, and evaluate cutting-edge AI research within the ecosystems of premier conferences directly advances these national priorities.

\subsubsection{Conclusion}

The totality of the evidence establishes that the beneficiary has achieved sustained national acclaim and is among the small percentage at the very top of the field of AI and machine learning systems development for critical infrastructure. The pattern of federal archival, cross-state government adoption, systemically important financial institution reliance, and selection to evaluate research within the ecosystems of premier AI/ML conferences demonstrates recognition that extends beyond any single employer or institution. The beneficiary's continued contributions to AI systems serving tens of millions of Americans and informing federal transportation policy constitute work of national importance that benefits the United States.

The final merits determination is satisfied.
